\addchap{\emph{A Note on How This Book Was Made}}
\draftstatus{Draft 1}

\lettrine{T}{here} is one more thing I owe the reader, and it is a plain account of how the object in your hands came to exist. I would rather give it than have somebody find it out, and I have a better reason than that, which I will come to at the end.

The ideas in this book are mine and they are old. The shadow push, the layers without a floor, charge as circulation, the copies, the local ratio: those were in my notes before I had any help writing them down, and the earliest of them were in my head before I could have defended them to anybody. The chronology in the epilogue is honest; nothing in it was supplied to me.

The sentences are another matter. Most of the prose in this book was drafted by an artificial intelligence working from my instructions, in a voice built deliberately from my own earlier writing, over a long conversation that ran to hundreds of exchanges. I set the arguments and the order; it proposed the words; I read every one of them and kept, corrected, or threw them out. Where it wrote something I would not have said, it was cut. Where it wrote something better than I could have, and true, it stayed --- and there are such places, and I am not going to point at them and blush.

It did other work, too, and the reader may as well know its extent. It drew most of the diagrams in code, so that their geometry would be arithmetic rather than eyeballing. The plates showing three-dimensional shapes go further: each states in advance what its own construction must satisfy --- that the grazing line really does meet the ball's radius at a right angle, that the blocked patch really is a quarter of its neighbour, that the small region marked in one panel really is the region the next panel magnifies --- and the program that made the figures in this book refuses to draw at all if one of those statements has stopped being true. I am fond of that arrangement, and not only because it caught mistakes. A picture that can fail its own test is the nearest thing a diagram can be to an honest witness. It ran the fact-checking: every number, date, page reference and attribution in this book was fetched against a named source, and the corrections came back in a report I read and ruled on. Several of them were corrections to my own errors. Several were corrections to \emph{its} errors, caught by procedures we had agreed on in advance for exactly that purpose. Two of the six wounds in Chapter~7 reached their final form because a machine told me a number I had used was wrong, and it was.

What it did not do was have the ideas, decide the physics, or supply the memories. It could not tell me whether my father explained light to me when I was eight. It could not decide whether the tower has a floor. When I told it that a collapsing star's descent need not be endless, it had the argument the wrong way round and said so; when I told it that a black hole's silhouette grows as the two-thirds power of its mass and not the square, it checked, agreed, and the sixth wound shrank to its true size. The judgements in this book are mine, and so is the responsibility for them.

Chapter~3 is a special case and I want to be exact about it. That dialogue is not a literary device. It is an edited record of an argument I actually had, with a machine, over many rounds, in which I set out to prove that a bottomless universe was the only possible kind and failed. The concessions in it run both ways, and mine are on the page because they are the most useful thing in the chapter. I lost that argument twice, and the book is better for it than the book I meant to write.

There is a record. The manuscript, the reference file, the fact-checking reports, and the whole trail of corrections with my approvals on them live in a repository that I intend to open to the public when this book is published. Anyone who wants to see the seams can see them.

And here is the better reason I mentioned. This book argues, for a hundred and some pages, that a thing is understood only when you can follow the chain of reasons under it, and that a wall with no explanation behind it is cheating wherever you find it. A book making that argument while concealing its own construction would be breaking its own rule on its last page. So: no magic here either. The ideas are mine, the labour was shared, the errors were caught by procedure, and you can go and look.

